\chapter{Studium przypadku}

% Odpowiednik sekcji Illustrative Examples. Pokazujesz, że narzędzie działa
% na danych, a nie tylko się kompiluje. Objętość: 8-10 stron.

Wprowadź zagadnienie, na którym demonstrujesz narzędzie, i uzasadnij jego wybór. Zagadnienie ma należeć do dziedziny wskazanej w tytule pracy.

\section{Charakterystyka danych}

Opisz zbiór: pochodzenie, licencję, liczbę obserwacji i zmiennych, skale pomiarowe, braki danych. Jeżeli używasz danych wygenerowanych własną procedurą, odeślij do sekcji, w której ją opisałeś, i podaj przyjęte parametry. Wyjaśnij, jak wybór zbioru wpływa na to, co wolno wywnioskować.

\section{Przebieg analizy}

Pokaż pełną ścieżkę od danych surowych do wyniku.

\begin{lstlisting}[caption={Przebieg analizy z uzyciem pakietu},label={lst:analiza}]
library(NazwaPakietu)

dane <- przygotuj_dane(zbior_surowy, model = "Kryterium =~ zmienna_a + zmienna_b")
wynik <- oblicz_wynik(dane, wagi = "entropia")

summary(wynik)
\end{lstlisting}

\section{Wyniki}

Przedstaw wyniki w tabeli i na rysunku, a następnie zinterpretuj je w tekście. Tabela i rysunek nie zastępują interpretacji.

\begin{rysunekepi}{Tytuł rysunku}{rys:wynik}
% \includegraphics[width=0.78\textwidth]{nazwa-pliku.png}
\zrodlo{oprac. własne}
\end{rysunekepi}

Oddziel to, co pokazuje narzędzie, od tego, co twierdzisz Ty. Narzędzie zwraca liczby; wniosek merytoryczny jest Twój i wymaga uzasadnienia.

\section{Porównanie i ocena poprawności}

Zestaw wyniki z rezultatem opublikowanym w artykule źródłowym albo z metodą odniesienia. Rozbieżność nie jest porażką – jest wynikiem, o ile potrafisz wskazać jej przyczynę.
